\section{An entropy bound on contacts}\label{sec:entropy}

Continue to assume coexistence in the coordinate direction, and use the two
tapes and their common blocks from Section~\ref{sec:bridges}.  We prove that
the contact count in \eqref{eq:contact-count} satisfies
\(m(J)=O(J/\log J)\).  The argument compares a portion of the actual tapes
with independent exact-cut bridges.  Contacts are unlikely for the independent
bridges, whereas the information needed to pass from those bridges to the
actual tapes can be bounded by the growth of a lateral-displacement entropy.
We adapt the comparison and contact argument of
\cite[Proposition~5.1 and Lemmas~5.2--5.3]{OpenAI2026}, using first proximity
arrivals to separate departure rows; all estimates needed here are proved
below.
All constants in this section may depend on the fixed environment law and
the fixed marking parameters, but not on the scales \(n\) or \(J\).

\subsection{Entropy of the lateral displacement}

For a countable random variable \(V\) with masses \(p(v)\), its Shannon
entropy \cite{Shannon1948} is
\[
 \Ent(V)=\sum_v p(v)\log\frac1{p(v)}.
\]
All logarithms in this section are natural.  For probability laws \(P,Q'\)
on the same countable set, their relative entropy
\cite{KullbackLeibler1951} is
\[
 D(P\Vert Q')=\sum_v P(v)\log\frac{P(v)}{Q'(v)}.
\]
A zero mass contributes zero, and a positive \(P\)-mass at a zero
\(Q'\)-mass gives infinite relative entropy.  Conditional entropy averages
the entropy of the conditional law.  Conditional mutual information is
defined by
\[
 I(U;V\mid C)
 =\sum_c P(C=c)
 D\bigl(P_{U,V\mid C=c}\,\Vert\,
          P_{U\mid C=c}\otimes P_{V\mid C=c}\bigr).
\]
We use these definitions even when \(U\) and \(V\) are entire finite path
lists, whose entropies need not be finite.

The log-sum inequality gives nonnegativity of relative entropy and its
decrease under taking a function of the variables.  Factoring joint masses
gives the chain rules.  In particular, comparison with a conditional product
law, while keeping the marginal of \(C\) fixed, gives
\begin{align}
 &D\bigl(P_{C,U,V}\,\Vert\,
          P_C Q'_{U\mid C}Q'_{V\mid C}\bigr)\notag\\
 &\qquad=I(U;V\mid C)
   +\mathbb E\,D(P_{U\mid C}\Vert Q'_{U\mid C})
   +\mathbb E\,D(P_{V\mid C}\Vert Q'_{V\mid C}).
 \label{eq:conditional-product-entropy}
\end{align}
These relative-entropy identities hold on countable spaces: one may first
use finite partitions and then take increasing limits, or factor the
conditional masses directly, since the negative parts of relative-entropy
sums are integrable.  When the variable whose entropy is being subtracted
has finite entropy, the same rules give the usual identity between an
entropy difference and mutual information.  Thus conditioning reduces
entropy in all the finite-entropy expressions below.

Put \(q=d-1\), and recall the iid lateral increments \(S_i'\in\ZZ^q\) of
the common blocks.  Define
\[
 h_j'=\Ent(S_1'+\cdots+S_j'),\qquad j\ge1.
\]
Equation~\eqref{eq:lateral-moment} gives \(\mathbf E|S_1'|<\infty\).
Consequently
\begin{equation}
 h_j'\le C+q\log(j+1).
 \label{eq:lateral-entropy-bound}
\end{equation}
Indeed, on \(\ZZ^q\) compare the displacement law with the probability
weights
\[
 r_j(x)=c_j^{-1}\exp\bigl(-|x|/(j+1)\bigr),
 \qquad
 c_j=\sum_{x\in\ZZ^q}\exp\bigl(-|x|/(j+1)\bigr)
       \le C(j+1)^q.
\]
The expected negative logarithm of \(r_j\) is bounded by
\(\log c_j+j\mathbf E|S_1'|/(j+1)\).  Nonnegativity of relative entropy,
or its finite-partition version followed by a limit, bounds the Shannon
entropy by this finite quantity.  This proves both finiteness and
\eqref{eq:lateral-entropy-bound}.  The sequence \(h_j'\) is nondecreasing:
conditioning \(S_1'+\cdots+S_{j+1}'\) on the independent last summand
gives entropy \(h_j'\), while removing that conditioning can only increase
entropy.

\subsection{Random chunks and their reference law}

Fix an integer \(n\ge1\).  Independently of the tapes, choose independent
integer counts with laws
\begin{equation}
 \begin{split}
 I_0,I_1,I_3&\text{ uniform on }\{n,\ldots,2n-1\},\\
 I_2&\text{ uniform on }\{7n,\ldots,8n-1\}.
 \end{split}
 \label{eq:chunk-counts}
\end{equation}
Discard a buffer of \(I_0\) common blocks, and record its lateral gap
\[
 d_0=\sum_{j=1}^{I_0}S_j'.
\]
Divide the next blocks into three consecutive chunks containing
\(I_1,I_2,I_3\) blocks.  For chunk \(i\), let \(A_i,B_i\) be its positive
and negative slab lists, both in downward tape order, and let \(H_i\) be
their common width.  Set
\[
 \mathcal A=(A_1,A_2,A_3),\qquad
 \mathcal B=(B_1,B_2,B_3),\qquad H=(H_1,H_2,H_3).
\]
The lists contain relative marked words; they do not include their spatial
placements.  The three chunks are independent of one another and of
\((I_0,d_0)\).  To see this, condition on the four counts: the chunks and
the buffer then use disjoint portions of an iid common-block sequence;
averaging over the independent counts preserves this product structure.

Let \(X_A,X_B\in\ZZ^q\) be the sums of the lateral displacements of all
words in \(\mathcal A,\mathcal B\), respectively, measured in the words'
original chronological directions.  Define
\[
 Z=d_0+X_A,\qquad T=I_0+I_1+I_2+I_3.
\]
Write \(P^{(n)}\) for the resulting law of
\((H,Z,\mathcal A,\mathcal B)\).  The widths, the buffer displacement,
and the chunk displacements have finite first moments by
Proposition~\ref{prop:common-blocks}.  The lattice comparison just used
therefore gives finite entropy for
\(H,d_0,X_A,X_B,Z\), with \(n\) fixed.

Each tuple \((H,Z,\mathcal A,\mathcal B)\) determines a pair of placed
paths as follows.  Put \(N=H_1+H_2+H_3\).  Start the positive path at
height zero and lateral position zero, and concatenate chunks \(3,2,1\),
reversing the order of each positive list.  Every word is still traversed
in its original chronological direction.  Start the negative path at
height \(N-1\) and lateral position \(Z\), and concatenate chunks
\(1,2,3\) in downward order.  Let \(E_{\rm mid}\) be the event that a
positive departure in chunk \(2\) is within \(\ell^\infty\)-distance
\(R\) of a negative departure in chunk \(2\).

Under \(P^{(n)}\), these are exactly the original three chunks after a
common translation.  At the top of the chunks their lateral gap was
\(d_0\), and the positive path gains displacement \(X_A\) from bottom to
top.  The negative starting coordinate relative to the positive bottom is
therefore \(d_0+X_A=Z\).  The one-level offset in their heights is the
offset in the tape construction.  Figure~\ref{fig:opposed-chunks} shows
this placement and the distinction between tape order and chronological
order.

\begin{figure}[t]
\centering
\begin{tikzpicture}[x=1cm,y=.86cm,>=Stealth,
 every node/.style={font=\small},
 posword/.style={draw=blue!60!black,fill=blue!6,line width=.6pt},
 negword/.style={draw=orange!65!black,fill=orange!7,line width=.6pt},
 guide/.style={densely dashed,gray!65,line width=.4pt}]
 \node[align=center] at (3.1,8.05)
   {Positive tape\\list order: $A_1,A_2,A_3$};
 \node[align=center] at (7.0,8.05)
   {Negative tape\\list order: $B_1,B_2,B_3$};
 \draw[posword] (2.2,.7) rectangle (4,2.7);
 \draw[posword,fill=blue!15] (2.2,2.7) rectangle (4,4.9);
 \draw[posword] (2.2,4.9) rectangle (4,6.7);
 \draw[negword] (6.1,.42) rectangle (7.9,2.42);
 \draw[negword,fill=orange!17] (6.1,2.42) rectangle (7.9,4.62);
 \draw[negword] (6.1,4.62) rectangle (7.9,6.42);
 \node at (3.1,1.7) {$A_3$};
 \node at (3.1,3.8) {$A_2$};
 \node at (3.1,5.8) {$A_1$};
 \node at (7,1.42) {$B_3$};
 \node at (7,3.52) {$B_2$};
 \node at (7,5.52) {$B_1$};
 \draw[->,very thick,blue!60!black] (4.35,1.05)--(4.35,6.32);
 \draw[->,very thick,orange!65!black] (5.75,6.07)--(5.75,.77);
 \node[rotate=90,font=\footnotesize,fill=white,inner sep=2pt]
   at (4.35,3.6) {chronological order};
 \node[rotate=90,font=\footnotesize,fill=white,inner sep=2pt]
   at (5.75,3.4) {chronological order};
 \foreach \yy/\ll in {.7/0,2.7/H_3,4.9/{N-H_1},6.7/N} {
   \draw[guide] (1.72,\yy)--(2.2,\yy);
   \node[anchor=east] at (1.64,\yy) {$\ll$};
 }
 \foreach \yy/\ll in {.42/{-1},2.42/{H_3-1},4.62/{N-H_1-1},6.42/{N-1}} {
   \draw[guide] (7.9,\yy)--(8.38,\yy);
   \node[anchor=west] at (8.46,\yy) {$\ll$};
 }
 \fill[blue!60!black] (3.1,.7) circle (1.7pt);
 \fill[blue!60!black] (3.1,6.7) circle (1.7pt);
 \fill[orange!65!black] (7,6.42) circle (1.7pt);
 \node[anchor=north,align=center,font=\footnotesize] at (3.1,.22)
   {positive start\\$(0,0)$};
 \node[anchor=north,align=center,font=\footnotesize] at (7,.02)
   {negative terminal};
 \node[anchor=south,font=\footnotesize,fill=white,inner sep=2pt]
   at (7,6.47) {start $(N-1,Z)$};
 \node[anchor=south,font=\footnotesize,fill=white,inner sep=2pt]
   at (3.1,6.75) {end $(N,X_A)$};
\end{tikzpicture}
\caption{The three chunks in the opposed arrangement, after translation
to put the positive starting site at $(0,0)$, with $N=H_1+H_2+H_3$.
Both tape lists are indexed from top to bottom. The positive path follows
chunks $3,2,1$, reversing each slab list while retaining the forward steps
of every word; the negative path follows chunks $1,2,3$ in their original
order. The one-level offset makes the departure heights in corresponding
chunks agree. Under the actual law $P^{(n)}$, the lateral gap at the top
is $Z-X_A=d_0$. The shaded middle chunks define $E_{\rm mid}$. Rectangles
show height ranges, not path traces; widths, heights, and lateral
placements are schematic.}
\label{fig:opposed-chunks}
\end{figure}


Define a second law \(Q^{(n)}\) on these same variables.  Keep the
\(P^{(n)}\)-marginal of \((H,Z)\); conditional on \((H,Z)\), sample the
six lists independently, with
\[
 A_i\sim\mathcal B_{H_i}^+,
 \qquad B_i\sim\mathcal B_{H_i}^-,
 \qquad i=1,2,3.
\]
These exact-cut bridge laws are defined by \eqref{eq:bridge-law}.
Since \(H_i\ge Mn\), their normalizing constants have the positive lower
bounds supplied by \eqref{eq:bridge-renewal}.
Thus the reference law keeps the three widths and the negative starting
position, and independently resamples the six relative lists.  The same
placement rule defines \(E_{\rm mid}\) under both laws.

\begin{proposition}[Entropy comparison]\label{prop:entropy-comparison}
For every integer \(n\ge1\), the actual chunk law and the reference law
above satisfy
\begin{equation}
 D(P^{(n)}\Vert Q^{(n)})\le C+h_{14n}'-h_n'.
 \label{eq:entropy-comparison}
\end{equation}
\end{proposition}

\begin{proof}
All entropies and expectations in this proof refer to the original
tape-and-count experiment.
First omit \(Z\), and let \(D_0\) denote the relative entropy between
the resulting marginals of \(P^{(n)}\) and \(Q^{(n)}\).  A pair of lists
of common width \(h\) determines the number of its common renewals,
counting the terminal width \(h\) and excluding zero, and hence determines
its parsing into common blocks.  If that number is in the count window
for chunk \(i\), the actual unconditioned probability of the list pair is
\[
 \frac1n
 \prod_{\gamma\in A_i}W(\gamma)
 \prod_{\gamma\in B_i}W(\gamma).
\]
If the number is outside the window, the probability is zero.  This follows
from the common-block construction and the slab weights: the specified
pair of lists determines exactly one allowed count.  In contrast, the
product of the two bridge laws at width \(h\) has the same product of
weights divided by \(\bar u_h^+\bar u_h^-\).  Thus, conditional on
\(H_i=h\), the density relative to the independent bridges is, on the
actual support, exactly
\begin{equation}
 \frac{\bar u_h^+\bar u_h^-}
      {nP^{(n)}(H_i=h)}.
 \label{eq:chunk-pair-density}
\end{equation}
Independence of the three actual chunks now gives
\begin{align}
 I(\mathcal A;\mathcal B\mid H)
 &\le D_0\notag\\
 &=\sum_{i=1}^3
   \left(\Ent(H_i)-\log n
          +\mathbb E\log(\bar u_{H_i}^+\bar u_{H_i}^-)\right)
 \le\sum_{i=1}^3(\Ent(H_i)-\log n)
 \le C.
 \label{eq:chunk-entropy-cost}
\end{align}
The information inequality follows from
\eqref{eq:conditional-product-entropy}.  For the last bound, use
\(\mathbb E H_i=\mathbf E K_1\,\mathbb E I_i=O(n)\) and compare the
law of \(H_i\) with weights proportional to \(e^{-h/n}\) on the
nonnegative integers.  The normalizer is \(O(n)\), so
\(\Ent(H_i)\le\log n+C\).

Reintroducing \(Z\) costs its conditional information with the lists:
\begin{equation}
 \begin{split}
 D(P^{(n)}\Vert Q^{(n)})
 &=D_0+I(Z;\mathcal A,\mathcal B\mid H)\\
 &=D_0+\Ent(Z\mid H)-\Ent(d_0).
 \end{split}
 \label{eq:buffer-information}
\end{equation}
Here the lists determine \(X_A\), and the buffer is independent of the
lists and widths; conditional on those data, \(Z\) is a translate of
\(d_0\).  Likewise, given \((\mathcal A,H)\), the variable \(Z\) is
conditionally independent of \(\mathcal B\).  Conditional data processing
and \eqref{eq:chunk-entropy-cost} imply
\[
 I(Z;\mathcal B\mid H)
 \le I(\mathcal A;\mathcal B\mid H)\le D_0.
\]
Every variable whose entropy is evaluated in the next calculation has
finite entropy; the conditioning lists need not:
\begin{align}
 \Ent(T,Z+X_B)
 &\ge\Ent(T,Z+X_B\mid\mathcal B,H)
   =\Ent(T,Z\mid\mathcal B,H)\notag\\
 &\ge\Ent(Z\mid H)-D_0
       +\Ent(T\mid Z,\mathcal A,\mathcal B,H)\notag\\
 &=\Ent(Z\mid H)-D_0+\Ent(I_0\mid d_0).
 \label{eq:buffer-entropy-lower}
\end{align}
For the final identity, each pair of chunk lists determines its common
block count, hence \(I_1,I_2,I_3\).  The lists and \(Z\) also determine
\(d_0=Z-X_A\).  Independence from the buffer then identifies the
remaining conditional entropy of \(T\) with \(\Ent(I_0\mid d_0)\).

Subtract \(\Ent(I_0,d_0)=\Ent(d_0)+\Ent(I_0\mid d_0)\) from
\eqref{eq:buffer-entropy-lower} and use
\eqref{eq:buffer-information}.  The result is
\begin{equation}
 D(P^{(n)}\Vert Q^{(n)})
 \le 2D_0+\Ent(T,Z+X_B)-\Ent(I_0,d_0).
 \label{eq:two-entropy-costs}
\end{equation}
Now \(Z+X_B\) is the sum of the lateral increments through the first
\(T\) common blocks.  Both \(T\) and \(I_0\) are independent of the
underlying block sequence.  It follows that
\begin{align*}
 \Ent(T,Z+X_B)-\Ent(I_0,d_0)
 &=\Ent(T)+\mathbb E h_T'-\log n-\mathbb E h_{I_0}'\\
 &\le C+h_{14n}'-h_n'.
\end{align*}
Indeed, \(T<14n\), \(I_0\ge n\), and \(T\) has at most \(4n\)
possible values.  Monotonicity of \(h_j'\) proves the displayed bound.
Finally, the term \(2D_0\) is bounded uniformly in \(n\) by
\eqref{eq:chunk-entropy-cost}, proving
\eqref{eq:entropy-comparison}.
\end{proof}

\subsection{Contacts under the reference law}

The entropy comparison will bound the contact count once we show that
\(E_{\rm mid}\) contains the relevant actual contacts and is unlikely
under the reference law.
For every realization of the counts, a contact in a positive common block
with index in \((4n,8n]\) implies \(E_{\rm mid}\).  In fact, chunk \(2\)
starts after at most \(4n-2\) common blocks and ends at or beyond block
\(9n\).  The negative departure realizing a contact is in the same or a
neighboring block, by the common-block geometry in
Section~\ref{sec:bridges}; all those blocks remain inside chunk \(2\).
Consequently
\begin{equation}
 \mathbf P\{\text{a contact in a positive block indexed in }(4n,8n]\}
 \le P^{(n)}(E_{\rm mid}).
 \label{eq:middle-contact-inclusion}
\end{equation}

\begin{proposition}[Reference contact bound]\label{prop:reference-contact}
For every integer \(n\ge1\),
\begin{equation}
 Q^{(n)}(E_{\rm mid})\le\delta_n,
 \qquad
 \delta_n=\min\{1,Cn(F_+(n)+F_-(n))\}.
 \label{eq:reference-contact}
\end{equation}
\end{proposition}

\begin{proof}
Fix \((H,Z)\) in the support of its retained marginal.  Conditional on
these data, all six reference lists are independent.  Reversing the
positive list order preserves its bridge law by
\eqref{eq:bridge-law}.  Retain the positive path formed by chunks \(3\)
and \(2\), and call it \(\alpha\).  It runs from the origin through its
first hit of height \(N-H_1\).  For the negative path formed by chunks
\(1\) and \(2\), retain only the prefix through its first arrival within
distance \(R\) of \(\Dep(\alpha)\); call the retained prefix \(\beta\)
and its terminal site \(y'\).  On \(E_{\rm mid}\) such an arrival exists
no later than a negative departure in chunk \(2\), and hence before the
two negative chunks terminate.
This first proximity need not itself join the two middle chunks; the
outer chunk widths will nevertheless put both visits at progress at
least \(n\).

The geometry puts this encounter away from both starting sides.  The two
negative chunks stay at or below \(N-1\) and end at their first hit of
\(H_3-1\).  Therefore
\[
 H_3\le y_1'\le N-1.
\]
Choose a departure \(y\) of \(\alpha\) within distance \(R\) of \(y'\).
Since \(\alpha\) ends on its first hit of \(N-H_1\),
\begin{equation}
 H_3-R\le y_1\le N-H_1-1,
 \qquad N-1-y_1'\ge H_1-R.
 \label{eq:reference-contact-heights}
\end{equation}
The inequalities \(H_1,H_3\ge Mn\) and \(M>R+1\) show that both visits
occur at or after the respective path has reached signed progress \(n\)
from its start.  Up to these visits neither path has hit either absolute
height \(-1\) or \(N\).

We next express the probability of these prefixes using two raw walks in
one environment.  The unused positive chunk \(1\) and negative chunk
\(3\) integrate to one.  The four remaining bridge normalizations,
\[
 \bar u_{H_3}^+\bar u_{H_2}^+
 \bar u_{H_1}^-\bar u_{H_2}^-,
\]
have a fixed positive lower bound by \eqref{eq:bridge-renewal}.  Within
each sign, concatenating the two bridges gives the product of their raw
weights, by the separation of active departures established in
Section~\ref{sec:bridges}.

For fixed \(\alpha\) and a retained negative prefix \(\beta\), sum over
all allowed negative completions to the terminal height \(H_3-1\).
Their total raw weight is at most \(W(\beta)\): full words to that height
end on its first hit, so the corresponding continuation events are
disjoint.  Nor is there multiplicity from the lists.  The intermediate
chunk boundary is a first hit of its prescribed height, and the bridge
parsing into slabs is unique.  The same uniqueness holds for the positive
concatenation.  Thus, after dropping any further completion restrictions,
the conditional contact probability is at most a fixed constant times a
sum of products
\begin{equation}
 W(\alpha)W(\beta)
 =\mathbb E_Q\bigl[p_\omega(\alpha)p_\omega(\beta)\bigr].
 \label{eq:reference-prefix-transfer}
\end{equation}
The equality holds for each retained pair: every departure of \(\beta\)
precedes its first proximity arrival and is therefore at distance strictly
greater than \(R\) from every departure of \(\alpha\).  The terminal
arrival \(y'\) uses no departure row.  Lemma~\ref{lem:separated-weights}
then applies, with repeated departures within either path retained in
their respective weights.

The right side of \eqref{eq:reference-prefix-transfer} is the probability
of the specified prefixes for two raw marked walks, started at \(0\) and
\((N-1,Z)\), independently conditional on a shared environment of law
\(Q\).  Summing these probabilities introduces no overcounting.  A full
positive trajectory has at most one prefix through its first hit of the
fixed height \(N-H_1\); that prefix determines \(\Dep(\alpha)\), and a
full negative trajectory has at most one prefix through its first arrival
within distance \(R\) of that set.

Only in this shared-environment experiment do we classify an encounter by
quenched exit probabilities.  For \(0\le x_1\le N-1\), define
\[
 q_x^\omega=P_{x,\omega}(T_{-1}<T_N),
\]
where \(T_a\) here is the first hit of absolute first-coordinate height
\(a\).  Slab exit is almost surely finite, so
\(1-q_x^\omega=P_{x,\omega}(T_N<T_{-1})\).  At the sites in
\eqref{eq:reference-contact-heights}, if \(q_y^\omega\ge1/2\), the
positive walk has visited, after reaching progress \(n\), a site with
probability at least \(1/2\) of exiting through its starting side's
barrier \(-1\).  If \(q_y^\omega<1/2\), join \(y'\) to \(y\) by a
coordinate path of at most \(dR\) steps inside the slab.  Uniform
ellipticity gives
\[
 1-q_{y'}^\omega
 \ge\kappa^{dR}(1-q_y^\omega)
 \ge\kappa^{dR}/2.
\]
In that case the negative walk has visited, after reaching its signed
progress \(n\), a site with probability at least \(\kappa^{dR}/2\) of
exiting through its own starting side's barrier \(N\).

For either sign \(\sigma\), consider the raw marginal event that, after
its first hit of signed progress \(n\) and before slab exit, the walk
visits a site whose quenched probability of exit through the starting
side's barrier is at least \(c=\kappa^{dR}/2\).  With the environment
fixed, stop at the first such visit.  The probability of the indicated
exit after this visit is at least \(c\).  Integrating the resulting
inequality over environments bounds the visit probability by
\begin{equation}
 c^{-1}(u_n^\sigma-u_N^\sigma).
 \label{eq:reference-visit-bound}
\end{equation}
Indeed, from each of the two starting sites, the barriers have signed
relative heights \(-1,N\).  The difference
\(u_n^\sigma-u_N^\sigma\) is exactly the raw probability of hitting
signed height \(n\) and then exiting at \(-1\) before \(N\): a path
reaching \(N\) must first reach \(n\), and exit from the finite slab is
almost surely finite.  Stationarity identifies the negative walk's
translated raw probabilities with \(u_n^-,u_N^-\).

Every retained pair belongs to at least one of these two raw visit events.
Combining \eqref{eq:reference-prefix-transfer}--\eqref{eq:reference-visit-bound}
therefore gives
\[
 Q^{(n)}(E_{\rm mid}\mid H,Z)
 \le C\sum_{\sigma\in\{+,-\}}(u_n^\sigma-u_N^\sigma)
 \le CN(F_+(n)+F_-(n)),
\]
where the last inequality is Lemma~\ref{lem:width-tail}.  This estimate
uses the raw marginal visit probabilities after the prefix transfer; it
does not condition an exit probability on the contact event.  Finally,
the retained marginal has
\[
 \mathbb E N=\mathbf E K_1\,\mathbb E(I_1+I_2+I_3)=O(n).
\]
Averaging over \((H,Z)\) and also using the trivial probability bound
one proves \eqref{eq:reference-contact}.
\end{proof}

\subsection{Summing over scales}

\begin{proposition}[Contact upper bound]\label{prop:contact-upper}
For the two independent opposed tapes under coordinate coexistence, the
contact count defined in \eqref{eq:contact-count} satisfies
\begin{equation}
 m(J)=O\!\left(\frac{J}{\log J}\right)
 \qquad(J\to\infty).
 \label{eq:contact-upper}
\end{equation}
\end{proposition}

\begin{proof}
Apply data processing in Proposition~\ref{prop:entropy-comparison} to
the indicator of \(E_{\rm mid}\).  If \(p=P^{(n)}(E_{\rm mid})\) and
\(q_n=Q^{(n)}(E_{\rm mid})\), binary relative entropy satisfies
\[
 D(\operatorname{Bern}(p)\Vert\operatorname{Bern}(q_n))
 \ge p\log(1/q_n)-\log2.
\]
Together with Proposition~\ref{prop:reference-contact}, this yields, when
\(\delta_n>0\),
\begin{equation}
 P^{(n)}(E_{\rm mid})\log(1/\delta_n)
 \le C+h_{14n}'-h_n'.
 \label{eq:contact-entropy-price}
\end{equation}
If \(\delta_n=0\), then \(q_n=0\), and the finite relative entropy in
\eqref{eq:entropy-comparison} forces
\(P^{(n)}(E_{\rm mid})=0\).  The same observation handles \(q_n=0\)
when \(\delta_n>0\).

For \(l\ge2\), take \(n=2^{l-2}\).  The width-tail summability in
Lemma~\ref{lem:width-tail} implies
\[
 \sum_{l\ge2}\delta_{2^{l-2}}<\infty.
\]
Among \(2\le l\le J\), at most \(O(\sqrt J)\) indices can have
\(\delta_{2^{l-2}}>J^{-1/2}\).  At every other index with positive
\(\delta_{2^{l-2}}\), the logarithm in
\eqref{eq:contact-entropy-price} is at least \(\tfrac12\log J\).
The corresponding entropy differences have a telescoping bound:
\begin{align*}
 \sum_{l=2}^J
   \bigl(h_{14\cdot2^{l-2}}'-h_{2^{l-2}}'\bigr)
 &\le\sum_{l=2}^J
   \bigl(h_{2^{l+2}}'-h_{2^{l-2}}'\bigr)\\
 &\le4h_{2^{J+2}}'=O(J),
\end{align*}
by monotonicity and \eqref{eq:lateral-entropy-bound}.  Since
\((4n,8n]=(2^l,2^{l+1}]\), use
\eqref{eq:middle-contact-inclusion}, sum
\eqref{eq:contact-entropy-price} over the remaining indices, and bound
each exceptional probability by one.  Adding the initial scale gives
\[
 m(J)\le 1+O(\sqrt J)+O(J/\log J)=O(J/\log J),
\]
as claimed.
\end{proof}
